\chapter{Introduction Générale}
\label{chap:introduction}

\section{Contexte du stage}

Dans le cadre de ma formation d'ingénieur en Informatique et Réseaux à l'EMSI --- École Marocaine des Sciences de l'Ingénieur, campus de Rabat, j'ai effectué un stage de cinq semaines du 2 juillet au 8 août 2026 au sein de la société \textbf{Axeli}, spécialisée dans les services informatiques et l'infogérance au Maroc.

Ce stage avait pour objectif principal de concevoir et de mettre en oeuvre un laboratoire de sécurité réseau complet, utilisant des pare-feux FortiGate dans un environnement virtualisé GNS3. L'expérience a permis de concilier les connaissances théoriques acquises en cours avec une mise en pratique professionnelle des technologies de sécurité réseau.

\section{Objectifs du stage}

Le stage s'est articulé autour de 16 objectifs définis par l'encadrant, couvrant l'ensemble des fonctionnalités d'un pare-feu FortiGate :

Les objectifs ont été regroupés en trois catégories selon leur état d'avancement :

\begin{table}[H]
\centering
\caption{Objectifs réalisés --- Fonctionnalités opérationnelles dans le laboratoire}
\label{tab:objectifs-realises}
\begin{tabularx}{\textwidth}{|c|X|c|}
\hline
\textbf{N°} & \textbf{Objectif} & \textbf{Statut} \\
\hline
1 & Configuration des firewall policies & Réalisé \\
2 & Configuration du NAT (SNAT) & Réalisé \\
3 & Routage statique & Réalisé \\
4 & Routage inter-LAN (Docker Router) & Réalisé \\
8 & High Availability (2 clusters actif-passif) & Réalisé \\
15 & Authentification locale & Réalisé \\
16 & Logging et supervision (mémoire + syslog) & Réalisé \\
\hline
\end{tabularx}
\end{table}

\begin{table}[H]
\centering
\caption{Objectifs partiellement réalisés --- Fonctionnalités configurées mais non testables en conditions réelles}
\label{tab:objectifs-partiels}
\begin{tabularx}{\textwidth}{|c|X|X|}
\hline
\textbf{N°} & \textbf{Objectif} & \textbf{Explication} \\
\hline
9 & Antivirus & Profils configurés sur FortiGate, signatures factory uniquement (pas de FortiGuard). L'endpoint EICAR existe côté Docker mais le moteur AV ne dispose pas de mises à jour. \\
\hline
10 & IPS & Mêmes limitations que l'AV. Les signatures factory couvrent les attaques les plus courantes mais sans mise à jour. \\
\hline
12 & Web Filtering & Listes d'URL statiques configurées. Sans FortiGuard, la catégorisation automatique n'est pas disponible. \\
\hline
13 & DNS Filtering & Filtres DNS statiques configurés. Mêmes limitations que le Web Filtering. \\
\hline
14 & SSL Deep Inspection & Certificat CA auto-signé généré. Non appliqué aux politiques firewall (limite de 3 politiques). \\
\hline
\end{tabularx}
\end{table}

\begin{table}[H]
\centering
\caption{Objectifs non réalisés --- Contraintes de licence et d'infrastructure}
\label{tab:objectifs-non-realises}
\begin{tabularx}{\textwidth}{|c|X|X|}
\hline
\textbf{N°} & \textbf{Objectif} & \textbf{Raison} \\
\hline
5 & Routage BGP & Remplacé par le Docker Router. Le routage dynamique n'a pas été implémenté faute de ports disponibles. \\
\hline
6 & VPN IPSec & Licence eval (chiffrements faibles) + firewall cloud (paquets bloqués avant réception). \\
\hline
7 & VPN SSL & Licence eval : 3 interfaces maximum, pas de port pour ssl.root. \\
\hline
11 & Application Control & Licence FortiGuard requise. Impossible sans abonnement. \\
\hline
\end{tabularx}
\end{table}

\noindent\textbf{Note :} Sur 16 objectifs, 7 ont été entièrement réalisés, 5 partiellement (profils configurés mais non testables sans FortiGuard), et 4 non réalisables (contraintes de licence ou d'infrastructure). La licence d'évaluation FortiGate a été le principal frein : elle limite l'appareil à 3 interfaces, 3 politiques firewall et 3 routes, et n'inclut pas les mises à jour de signatures FortiGuard.

\section{Problématique}

Comment concevoir un environnement de sécurité réseau fonctionnel et reproductible en utilisant des équipements virtualisés, tout en tenant compte des contraintes de licence et de matériel inhérentes à un environnement d'évaluation ? Cette problématique a guidé l'ensemble des choix d'architecture, de la topologie réseau jusqu'aux solutions de contournement mises en oeuvre pour chaque fonctionnalité indisponible.
